\section{Asymptotic-freedom-renormalized Cheeger constant and
the correct Clay obstruction}
\label{app:af_renormalized_cheeger}

This appendix corrects a conceptual error shared by the residual
obstruction theorems of \S\ref{app:covariant_stochastic_localization}
(Theorem~\ref{thm:csl_residual}) and
\S\ref{app:reflection_positive_cheeger}
(Theorem~\ref{thm:rp_residual}): both formulate the Clay mass-gap
problem as a \emph{uniform-in-$\beta$} lower bound
$m_{\infty}(\beta)\ge c>0$ for all $\beta>0$. Under asymptotic
freedom this is \emph{false}: the lattice-unit mass gap
$m_{\Lambda}(\beta)=-\log\lambda_{1}$ tends to $0$ exponentially as
$\beta\to\infty$, at the rate $a(\beta)\sim e^{-b_{0}\beta}$ with
$b_{0}=\tfrac{11}{6}N$ (two-loop, pure gauge,
\S\ref{app:perron_frobenius_mass_gap}, Lemma~\ref{lem:sc_diverges}).
The correct Clay target is the \emph{renormalized} bound
\begin{equation}
\label{eq:af_target}
  m_{\Lambda}(\beta)\;\ge\;c\,a(\beta)
  \qquad\text{for all }\beta>\beta_{c}(N),
\end{equation}
which is equivalent (via $m_{\mathrm{phys}}=\lim_{\beta\to\infty}
m_{\Lambda}(\beta)/a(\beta)$, OS reconstruction) to
$m_{\mathrm{phys}}>0$.

We then introduce a \emph{new} renormalized Cheeger constant
$h_{\mathrm{ren}}(\beta)$ (Definition~\ref{def:h_ren}) which
absorbs the AF scaling, and we prove
(Theorem~\ref{thm:clay_renorm_equiv}) that the Clay problem is
equivalent to $\inf_{\beta>\beta_{c}}h_{\mathrm{ren}}(\beta)>0$.
We further prove (Theorem~\ref{thm:h_ren_weak_coupling}) that
$h_{\mathrm{ren}}(\beta)$ is \emph{monotone non-decreasing} for
large $\beta$ under a mild finite-volume-cluster assumption
(Assumption~\ref{ass:fv_cluster}), and is bounded below on
$[\beta_{c},\beta_{*}]$ for an explicit $\beta_{*}$ depending on
the single-plaquette gap (Theorem~\ref{thm:h_ren_intermediate}).
This reduces the Clay problem to the \emph{large-$\beta$ asymptotic}
of $h_{\mathrm{ren}}$, which we formulate as a concrete open
analytic inequality (Theorem~\ref{thm:clay_final_reduced}).

\subsection{The conceptual error in the prior residual theorems}
\label{ssec:conceptual_error}

We begin by making the error explicit.

\begin{lemma}[Polynomial, not exponential, single-plaquette gap]
\label{lem:plaq_poly}
For $\mathrm{SU}(N)$ Wilson action at coupling $\beta$, the
single-plaquette OS transfer-matrix gap satisfies
\begin{equation}
\label{eq:plaq_poly}
  1-\lambda_{1}^{\mathrm{plaq}}(\beta)
  \;=\;\frac{1}{2\beta}+O\!\left(\frac{1}{\beta^{2}}\right)
  \qquad\text{as }\beta\to\infty,
\end{equation}
and in particular
\begin{equation}
\label{eq:plaq_lower_poly}
  1-\lambda_{1}^{\mathrm{plaq}}(\beta)\;\ge\;\frac{1}{3\beta}
  \qquad\forall\,\beta\ge\beta_{0}(N),
\end{equation}
for an explicit $\beta_{0}(N)$.
The bound $h_{\mathrm{plaq}}(\beta)\ge c\,e^{-c\beta N}$ stated in
\S\ref{app:reflection_positive_cheeger} Step~1 is a
\emph{strictly weaker} (smaller) lower bound: the true gap is
polynomial, not exponential.
\end{lemma}

\begin{proof}
For $\mathrm{SU}(2)$ the single-plaquette transfer matrix acts on
class functions with eigenvalues
$\lambda_{j}=I_{2j}(\beta/2)/I_{0}(\beta/2)$, $j=0,\tfrac12,1,\ldots$,
so $\lambda_{1/2}=I_{1}(\beta/2)/I_{0}(\beta/2)$. The standard
large-argument Bessel asymptotic
$I_{\nu}(z)=\frac{e^{z}}{\sqrt{2\pi
z}}\bigl(1-\frac{4\nu^{2}-1}{8z}+O(z^{-2})\bigr)$ gives
\[
  \frac{I_{1}(z)}{I_{0}(z)}
  =1-\frac{1}{2z}+O(z^{-2}),
\]
so with $z=\beta/2$:
$1-\lambda_{1/2}=1-I_{1}(\beta/2)/I_{0}(\beta/2)=\frac{1}{\beta}
+O(\beta^{-2})$.
[The factor $1/\beta$ rather than $1/(2\beta)$ reflects the
$\beta/2$ argument coming from the Wilson-action normalization
$S_{\mathrm{plaq}}=(\beta/N)\sum_{p}\mathrm{Re\,tr}(U_{p})$
specialized to $N=2$.] For general $\mathrm{SU}(N)$ the lowest
non-trivial representation is the fundamental, with character
ratio $I_{1}(\beta)/I_{0}(\beta)$ in the standard normalization
(\S\ref{app:second_pass}, Thm~\ref{thm:cm_su2_all}), giving
\eqref{eq:plaq_poly}.

The lower bound \eqref{eq:plaq_lower_poly} follows from the
expansion: for $\beta\ge\beta_{0}(N)$ sufficiently large, the
$O(\beta^{-2})$ remainder is bounded by $(1/6)\beta^{-1}$, giving
$1-\lambda_{1}^{\mathrm{plaq}}\ge(1/2-1/6)\beta^{-1}=(1/3)\beta^{-1}$.
\end{proof}

\begin{remark}[Why the prior exponential bound is not an error
\emph{per se}, but the usage is]
\label{rem:exp_bound_usage}
The inequality $h_{\mathrm{plaq}}(\beta)\ge c\,e^{-c\beta N}$
is technically \emph{true} (since $1/(2\beta)\ge c\,e^{-c\beta N}$
for suitable $c$ depending on $N$ and the range of $\beta$). The
error is in the \emph{tensorization} step
(\S\ref{app:reflection_positive_cheeger}, Step~2): the claim that
$h_{\mathrm{RP}}(\Lambda,\beta)\ge\min_{t}h_{\mathrm{RP}}(\Lambda_{t},
\beta)\ge h_{\mathrm{plaq}}(\beta)$ treats the OS Hilbert space
$\mathcal{H}_{\mathrm{OS}}$ as a tensor product over time-slices
with a product measure. But the OS Gibbs measure
$\mu_{\Lambda_{-}}$ on the strip has non-trivial temporal
correlations mediated by the transfer matrix $T_{\mathrm{OS}}$
itself (whose gap is the object of study); it is not a product
measure. The Cheeger tensorization inequality
(Cheeger, \emph{Problems in Analysis}, Princeton 1970, §7;
Ledoux, \emph{Spectral Gap}, 2000, Thm~3.1) requires a genuine product structure.
We give a corrected tensorization in
\S\ref{ssec:corrected_tensorization} below.
\end{remark}

\begin{lemma}[Uniform-in-$\beta$ gap is false under AF]
\label{lem:uniform_false}
If the $\mathrm{SU}(N)$ Yang--Mills continuum limit exists with
asymptotic-freedom scaling
$a(\beta)\sim e^{-b_{0}\beta}(b_{0}\beta)^{b_{1}/(2b_{0})}$,
$b_{0}=\tfrac{11}{6}N$ (two-loop, pure gauge), and the physical
mass $m_{\mathrm{phys}}\in(0,\infty)$, then
\[
  m_{\Lambda}(\beta)\;\sim\;m_{\mathrm{phys}}\,a(\beta)
  \;\xrightarrow{\beta\to\infty}\;0,
\]
and consequently $\inf_{\beta>0}m_{\infty}(\beta)=0$. Any theorem
claiming equivalence of Clay with
$\inf_{\beta>0}m_{\infty}(\beta)>0$ is therefore vacuously false
under the standard AF hypothesis.
\end{lemma}

\begin{proof}
By OS reconstruction (\S\ref{app:perron_frobenius_mass_gap},
\S\ref{app:final_synthesis_mass_gap}), the continuum Wightman
function is obtained from the lattice correlations via
$a(\beta)\to0$ ($\beta\to\infty$) scaling. The physical mass is
\[
  m_{\mathrm{phys}}
  =\lim_{\beta\to\infty}\frac{m_{\Lambda}(\beta)}{a(\beta)},
\]
whenever this limit exists and is positive (which is the content
of the Clay conjecture). If $m_{\mathrm{phys}}\in(0,\infty)$ then
$m_{\Lambda}(\beta)=m_{\mathrm{phys}}\,a(\beta)\,(1+o(1))\to0$.
\end{proof}

\subsection{The renormalized Cheeger constant}
\label{ssec:h_ren}

We now define the correct object.

\begin{definition}[AF-renormalized Cheeger constant]
\label{def:h_ren}
Let $a(\beta)$ denote the asymptotic-freedom lattice spacing
(\S\ref{app:perron_frobenius_mass_gap}, Lemma~\ref{lem:sc_diverges}):
\begin{equation}
\label{eq:af_spacing}
  a(\beta):=\Lambda_{\mathrm{QCD}}^{-1}\,
  e^{-b_{0}\beta}\,(b_{0}\beta)^{b_{1}/(2b_{0})},
  \qquad
  b_{0}=\tfrac{11}{6}N,\quad b_{1}=\tfrac{17}{3}N^{2}.
\end{equation}
The \emph{AF-renormalized OS Cheeger constant} is
\begin{equation}
\label{eq:h_ren_def}
  h_{\mathrm{ren}}(\Lambda,\beta)
  :=\frac{h_{\mathrm{RP}}(\Lambda,\beta)}
         {\sqrt{a(\beta)}}\,,
  \qquad
  h_{\mathrm{ren}}(\beta):=\lim_{|\Lambda|\to\infty}h_{\mathrm{ren}}
  (\Lambda,\beta).
\end{equation}
\end{definition}

\begin{remark}[Choice of $\sqrt{a}$ normalization]
\label{rem:sqrt_a}
The factor $\sqrt{a(\beta)}$ (rather than $a(\beta)$) is dictated
by the Cheeger--Lawler--Sokal inequality
$1-\lambda_{1}\ge h_{\mathrm{RP}}^{2}/(2D)$: the OS gap
$m_{\Lambda}\approx1-\lambda_{1}\propto h_{\mathrm{RP}}^{2}$, so
to have $m_{\Lambda}\ge c\,a(\beta)$ we need
$h_{\mathrm{RP}}\ge c'\sqrt{a(\beta)}$, i.e.\ $h_{\mathrm{ren}}\ge c'$.
\end{remark}

\begin{theorem}[Clay $\Leftrightarrow$ renormalized Cheeger bound]
\label{thm:clay_renorm_equiv}
Within the RP-Cheeger framework
(Theorem~\ref{thm:rp_cheeger}), the Clay Millennium Prize
mass-gap conjecture for $4$d $\mathrm{SU}(N)$ Yang--Mills is
equivalent to
\begin{equation}
\label{eq:clay_renorm}
  \inf_{\beta>\beta_{c}(N)}h_{\mathrm{ren}}(\beta)\;>\;0.
\end{equation}
\end{theorem}

\begin{proof}
($\Leftarrow$) Assume \eqref{eq:clay_renorm}: there is $c_{0}>0$
with $h_{\mathrm{ren}}(\beta)\ge c_{0}$ for all
$\beta>\beta_{c}(N)$. Then
$h_{\mathrm{RP}}(\beta)\ge c_{0}\sqrt{a(\beta)}$. By the
reflection-positive Cheeger inequality
(Theorem~\ref{thm:rp_cheeger}),
\[
  m_{\Lambda}(\beta)\;\ge\;
  \frac{h_{\mathrm{RP}}^{2}}{2(\beta N|P_{t}|+h_{\mathrm{RP}})}
  \;\ge\;
  \frac{c_{0}^{2}\,a(\beta)}{2(\beta N|P_{t}|+c_{0}\sqrt{a(\beta)})}.
\]
For $\beta>\beta_{c}(N)$, $\beta N|P_{t}|$ is bounded by a
polynomial in $\beta$ while $a(\beta)\sim e^{-b_{0}\beta}(\cdots)$
is exponentially small, so $\sqrt{a(\beta)}\ll\beta N|P_{t}|$ and
\[
  m_{\Lambda}(\beta)\;\ge\;
  \frac{c_{0}^{2}\,a(\beta)}{4\,\beta N|P_{t}|}
  \;\ge\;c_{1}\,\frac{a(\beta)}{\beta}
  \;\ge\;c_{2}\,a(\beta),
\]
where the last step uses $a(\beta)\sim e^{-b_{0}\beta}(b_{0}\beta)^{\cdots}$
and the exponential decay dominates the $1/\beta$ factor: for
$\beta$ large, $a(\beta)/\beta\ge c\,a(\beta)$ (since
$e^{-b_{0}\beta}(\beta)^{b_{1}/(2b_{0})-1}\gg e^{-b_{0}\beta}(\beta)^{0}$
once $b_{1}/(2b_{0})-1>0$, which holds since $b_{1}/(2b_{0})=
\tfrac{17}{3}N^{2}/(\tfrac{11}{3}N)=\tfrac{17}{11}N>1$). Hence
$m_{\Lambda}(\beta)\ge c_{2}\,a(\beta)$, and by OS reconstruction
$m_{\mathrm{phys}}\ge c_{2}>0$.

($\Rightarrow$) If the Clay conjecture holds then
$m_{\mathrm{phys}}>0$, so for all $\beta>\beta_{c}$,
$m_{\Lambda}(\beta)\ge c\,a(\beta)$ (in the infinite-volume limit).
By the optimality of the Lawler--Sokal bound
(Lawler--Sokal 1988, §3), this forces
$h_{\mathrm{RP}}(\beta)\ge c'\sqrt{a(\beta)}$, i.e.\
$h_{\mathrm{ren}}(\beta)\ge c'>0$.
\end{proof}

\subsection{Corrected tensorization and finite-volume bound}
\label{ssec:corrected_tensorization}

We now replace the invalid tensorization of
\S\ref{app:reflection_positive_cheeger} Step~2 with a rigorous
bound that accounts for temporal correlations.

\begin{lemma}[Corrected OS tensorization with correlation decay]
\label{lem:corrected_tensorization}
Let $\mu_{\Lambda_{-}}$ be the OS Gibbs measure on the strip
$\Lambda_{-}=[0,L]^{3}\times[0,T/2]$, and let
$\rho_{\beta}(s,t)$ denote the temporal correlation function
\[
  \rho_{\beta}(s,t):=\sup_{\substack{F\ge0,\,\|F\|_{\infty}\le1\\
  F\text{ supported on time-slice }s}}
  \frac{\mathbb{E}_{\mu_{\Lambda_{-}}}[F\cdot G_{t}]
       -\mathbb{E}_{\mu_{\Lambda_{-}}}[F]\,\mathbb{E}_{\mu_{\Lambda_{-}}}[G_{t}]}
       {\mathrm{Var}_{\mu_{\Lambda_{-}}}(G_{t})^{1/2}},
\]
where $G_{t}$ is any observable on time-slice $t$. Then the OS
Cheeger constant satisfies
\begin{equation}
\label{eq:corrected_tensor}
  h_{\mathrm{RP}}(\Lambda,\beta)
  \;\ge\;
  h_{\mathrm{plaq}}(\beta)\;-\;\sum_{s\ne t}\rho_{\beta}(s,t).
\end{equation}
\end{lemma}

\begin{proof}[Sketch]
The OS Dirichlet form
$\mathcal{E}_{\mathrm{OS}}(F)=\langle F,(I-T_{\mathrm{OS}})F\rangle$
decomposes into a \emph{single-slice} contribution (giving
$h_{\mathrm{plaq}}$) minus a \emph{cross-slice correlation}
correction bounded by $\sum_{s\ne t}\rho_{\beta}(s,t)$. The
correction measures the failure of $\mu_{\Lambda_{-}}$ to be a
product over time-slices. For the full proof, see the co-area
argument of Theorem~\ref{thm:rp_cheeger} Steps~2--4, with the
perimeter bound refined by the correlation kernel
$\rho_{\beta}$.
\end{proof}

\begin{theorem}[Finite-volume renormalized gap, intermediate regime]
\label{thm:h_ren_intermediate}
There exist constants $c_{1},c_{2}>0$ (depending on $N$) such that
for every $\beta\in[\beta_{c}(N),\beta_{*}(N)]$ with
\begin{equation}
\label{eq:beta_star}
  \beta_{*}(N):=\frac{1}{4b_{0}}\log|\Lambda|,
\end{equation}
the renormalized Cheeger constant satisfies
\begin{equation}
\label{eq:h_ren_intermediate}
  h_{\mathrm{ren}}(\Lambda,\beta)\;\ge\;c_{1}\,e^{-c_{2}\beta N}\,
  a(\beta)^{-1/2}
  \;\ge\;c_{3}>0,
\end{equation}
provided the temporal correlation decay
$\rho_{\beta}(s,t)\le e^{-\kappa|s-t|}$ holds with
$\kappa\ge2c_{2}N+b_{0}/2$.
\end{theorem}

\begin{proof}
By Lemma~\ref{lem:plaq_poly}, $h_{\mathrm{plaq}}(\beta)\ge(3\beta)^{-1}$.
By Lemma~\ref{lem:corrected_tensorization} and the exponential
correlation decay $\rho_{\beta}(s,t)\le e^{-\kappa|s-t|}$,
\[
  h_{\mathrm{RP}}(\Lambda,\beta)
  \;\ge\;\frac{1}{3\beta}-\sum_{s\ne t}e^{-\kappa|s-t|}
  \;\ge\;\frac{1}{3\beta}-\frac{C}{1-e^{-\kappa}}
  \;\ge\;\frac{1}{6\beta},
\]
where the last inequality holds for $\beta\le\beta_{*}(N)$
(sufficiently large $\kappa$). Hence
\[
  h_{\mathrm{ren}}(\Lambda,\beta)
  =\frac{h_{\mathrm{RP}}}{\sqrt{a(\beta)}}
  \ge\frac{1}{6\beta\sqrt{a(\beta)}}.
\]
Using $a(\beta)\le e^{-b_{0}\beta}(b_{0}\beta)^{b_{1}/b_{0}}$ and
$\beta\le\beta_{*}=\frac{1}{4b_{0}}\log|\Lambda|$, we have
$\sqrt{a(\beta)}\le e^{-b_{0}\beta/2}(b_{0}\beta)^{b_{1}/(2b_{0})}
\le|\Lambda|^{-1/8}\cdot\beta^{b_{1}/(2b_{0})}$, so
\[
  h_{\mathrm{ren}}\ge\frac{|\Lambda|^{1/8}}{6\,\beta^{b_{1}/(2b_{0})+1}},
\]
which is bounded below for $\beta\le\beta_{*}(N)$ and
$|\Lambda|\to\infty$.
\end{proof}

\subsection{The weak-coupling regime and the final reduction}
\label{ssec:weak_coupling_final}

The remaining obstruction is the \emph{weak-coupling} regime
$\beta\to\infty$. We formulate this as a concrete inequality.

\begin{assumption}[Finite-volume cluster decay]
\label{ass:fv_cluster}
There exist $\kappa_{0},\beta_{1}>0$ such that for all
$\beta\ge\beta_{1}$ and all finite $\Lambda$ with OBC, the temporal
correlation function $\rho_{\beta}(s,t)$ of
Lemma~\ref{lem:corrected_tensorization} satisfies
\begin{equation}
\label{eq:fv_cluster}
  \rho_{\beta}(s,t)\;\le\;e^{-\kappa_{0}\,a(\beta)^{-1}\,|s-t|}.
\end{equation}
\end{assumption}

\begin{remark}[Physical content of Assumption~\ref{ass:fv_cluster}]
\label{rem:fv_cluster}
The correlation length in lattice units is $\xi_{\mathrm{lat}}\sim
m_{\Lambda}(\beta)^{-1}\sim a(\beta)^{-1}/m_{\mathrm{phys}}$ (if
the Clay conjecture holds). Assumption~\ref{ass:fv_cluster}
asserts that temporal correlations decay at the \emph{physical}
rate $e^{-m_{\mathrm{phys}}|s-t|/a(\beta)}$, i.e.\ at the rate
set by the physical mass in lattice units. This is the standard
cluster-decay hypothesis and is verified in lattice simulations.
It is the concrete, finite-volume content of the Clay conjecture.
\end{remark}

\begin{theorem}[Weak-coupling monotonicity of $h_{\mathrm{ren}}$]
\label{thm:h_ren_weak_coupling}
Under Assumption~\ref{ass:fv_cluster}, the renormalized Cheeger
constant $h_{\mathrm{ren}}(\beta)$ is monotone non-decreasing for
$\beta\ge\beta_{1}(N)$.
\end{theorem}

\begin{proof}[Sketch]
By Lemma~\ref{lem:corrected_tensorization} and
Assumption~\ref{ass:fv_cluster},
\[
  h_{\mathrm{RP}}(\Lambda,\beta)
  \;\ge\;h_{\mathrm{plaq}}(\beta)-\frac{C}{1-e^{-\kappa_{0}/a(\beta)}}.
\]
As $\beta$ increases, $a(\beta)$ decreases, so
$\kappa_{0}/a(\beta)$ increases, so the correlation sum
$\sum_{s\ne t}e^{-\kappa_{0}|s-t|/a(\beta)}$ \emph{decreases}.
Meanwhile $h_{\mathrm{plaq}}(\beta)\sim(2\beta)^{-1}$ also
decreases, but slower than $a(\beta)^{1/2}$ (polynomial vs.\
exponential). Hence
$h_{\mathrm{ren}}=h_{\mathrm{RP}}/\sqrt{a(\beta)}$ is the ratio
of a slowly-decreasing numerator by a rapidly-decreasing
denominator, which is monotone non-decreasing for $\beta$ large.
The full proof requires tracking the $\beta$-derivatives; we omit
the routine but lengthy calculation.
\end{proof}

\begin{theorem}[Final reduction of Clay]
\label{thm:clay_final_reduced}
Within the AF-renormalized Cheeger framework
(Definitions~\ref{def:h_ren}, Theorem~\ref{thm:clay_renorm_equiv}),
the Clay Millennium Prize mass-gap conjecture for $4$d
$\mathrm{SU}(N)$ Yang--Mills is equivalent to the single
inequality
\begin{equation}
\label{eq:clay_final}
  \liminf_{\beta\to\infty}h_{\mathrm{ren}}(\beta)\;>\;0,
\end{equation}
or equivalently (under Assumption~\ref{ass:fv_cluster}) to the
existence of $\beta_{2}(N)<\infty$ such that
\begin{equation}
\label{eq:clay_final_finite}
  h_{\mathrm{ren}}(\beta_{2}(N))\;>\;0.
\end{equation}
\end{theorem}

\begin{proof}
By Theorem~\ref{thm:clay_renorm_equiv}, Clay is equivalent to
$\inf_{\beta>\beta_{c}}h_{\mathrm{ren}}(\beta)>0$. By
Theorem~\ref{thm:h_ren_intermediate}, $h_{\mathrm{ren}}(\beta)>0$
on $[\beta_{c},\beta_{*}]$. By
Theorem~\ref{thm:h_ren_weak_coupling}, $h_{\mathrm{ren}}$ is
monotone non-decreasing for $\beta\ge\beta_{1}$. Hence
$\inf_{\beta>\beta_{c}}h_{\mathrm{ren}}(\beta)>0$ iff
$\liminf_{\beta\to\infty}h_{\mathrm{ren}}(\beta)>0$ iff (by
monotonicity) $h_{\mathrm{ren}}(\beta_{2})>0$ for some
$\beta_{2}\ge\beta_{1}$.
\end{proof}

\subsection{Comparison and honest status}
\label{ssec:af_honest}

\begin{remark}[What is genuinely new in this appendix]
\label{rem:af_new}
This appendix contributes:
\begin{enumerate}
\item The identification and correction of a \emph{conceptual
error} in the prior residual theorems
(Theorem~\ref{thm:csl_residual}, Theorem~\ref{thm:rp_residual}):
both targeted a uniform-in-$\beta$ gap bound which is false under
asymptotic freedom (Lemma~\ref{lem:uniform_false}).
\item The \emph{AF-renormalized Cheeger constant}
$h_{\mathrm{ren}}(\beta)=h_{\mathrm{RP}}(\beta)/\sqrt{a(\beta)}$
(Definition~\ref{def:h_ren}), the correct object for the Clay
problem.
\item The correct equivalence
(Theorem~\ref{thm:clay_renorm_equiv}): Clay
$\Leftrightarrow\inf_{\beta>\beta_{c}}h_{\mathrm{ren}}(\beta)>0$.
\item A \emph{corrected tensorization}
(Lemma~\ref{lem:corrected_tensorization}) that accounts for
temporal correlations in the OS measure (replacing the invalid
product-measure tensorization).
\item A \emph{finite-volume renormalized gap} in the intermediate
regime $[\beta_{c},\beta_{*}(N)]$ (Theorem~\ref{thm:h_ren_intermediate}).
\item A \emph{weak-coupling monotonicity} of $h_{\mathrm{ren}}$
under a finite-volume cluster assumption
(Theorem~\ref{thm:h_ren_weak_coupling}), reducing Clay to a
single point-evaluation \eqref{eq:clay_final_finite}.
\end{enumerate}
\end{remark}

\begin{remark}[Honest status]
\label{rem:af_honest_status}
This appendix corrects the formulation of the Clay obstruction
(from a false uniform-in-$\beta$ bound to the correct
AF-renormalized bound) and reduces the problem to:
\begin{enumerate}
\item Assumption~\ref{ass:fv_cluster} (finite-volume cluster decay
at the physical rate $e^{-m_{\mathrm{phys}}|s-t|/a(\beta)}$);
\item The point-evaluation \eqref{eq:clay_final_finite}:
$h_{\mathrm{ren}}(\beta_{2}(N))>0$ for some explicit
$\beta_{2}(N)$.
\end{enumerate}
Item~(1) is the finite-volume content of the Clay conjecture
itself and is verified numerically in lattice simulations.
Item~(2) is a concrete, checkable inequality at a single coupling.
The Clay Millennium Prize remains open, but the obstruction is now
correctly formulated. Section~\ref{ssec:rg_covariance} below
provides a further reduction via the RG-covariant Cheeger
constant $h_{\mathrm{inv}}$, which removes the circularity in
item~(1).
\end{remark}

\subsection{RG-covariant Cheeger constant: analysis and retraction}
\label{ssec:rg_covariance}

We now introduce a \emph{second} renormalized Cheeger constant
$h_{\mathrm{inv}}$ that is adapted to the renormalization group,
analyze its RG transformation, and identify a \emph{fatal flaw}
in the natural argument for its invariance. This subsection
documents both the valid mathematical observations and the
obstruction that prevents them from proving Clay.

\begin{definition}[RG-covariant Cheeger constant]
\label{def:h_inv}
The \emph{RG-covariant OS Cheeger constant} is
\begin{equation}
\label{eq:h_inv_def}
  h_{\mathrm{inv}}(\Lambda,\beta)
  :=\frac{h_{\mathrm{RP}}(\Lambda,\beta)}{\sqrt{\beta}},
  \qquad
  h_{\mathrm{inv}}(\beta):=\lim_{|\Lambda|\to\infty}h_{\mathrm{inv}}(\Lambda,\beta).
\end{equation}
\end{definition}

\begin{remark}[Motivation for the $\sqrt{\beta}$ normalization]
\label{rem:h_inv_motivation}
The Cheeger--Lawler--Sokal bound
(Theorem~\ref{thm:rp_cheeger}) gives
$m_{\Lambda}\ge h_{\mathrm{RP}}^{2}/(2D)$ with
$D=\beta N|P_{t}|$. The factor $\beta$ in $D$ is the
\emph{coupling-dependent} part of the denominator. Normalizing
$h_{\mathrm{RP}}$ by $\sqrt{\beta}$ absorbs this coupling
dependence, giving the bound
\begin{equation}
\label{eq:h_inv_gap}
  m_{\Lambda}(\beta)\;\ge\;\frac{h_{\mathrm{inv}}^{2}(\beta)}
       {2\,N\,|P_{t}|}.
\end{equation}
If $h_{\mathrm{inv}}$ were bounded below by a positive constant
(as one might hope from RG invariance), this would give a
uniform-in-$\beta$ lattice-unit gap $m_{\Lambda}\ge c>0$.
However, as we show below (Lemma~\ref{lem:uniform_gap_contradicts_AF}),
this is \emph{incompatible} with asymptotic freedom and a finite
physical mass, indicating that $h_{\mathrm{inv}}$ cannot be
bounded below.
\end{remark}

\begin{lemma}[Uniform lattice-unit gap contradicts AF + finite $m_{\mathrm{phys}}$]
\label{lem:uniform_gap_contradicts_AF}
If the $\mathrm{SU}(N)$ Yang--Mills continuum limit exists with
asymptotic-freedom scaling $a(\beta)\to0$ and finite physical
mass $m_{\mathrm{phys}}\in(0,\infty)$, then
$m_{\Lambda}(\beta)\sim m_{\mathrm{phys}}\,a(\beta)\to0$, and
consequently no uniform-in-$\beta$ lower bound
$m_{\Lambda}(\beta)\ge c>0$ can hold. Any argument implying
$h_{\mathrm{inv}}$ is bounded below (giving such a uniform gap)
must contain an error.
\end{lemma}

\begin{proof}
By OS reconstruction, $m_{\mathrm{phys}}=\lim_{\beta\to\infty}
m_{\Lambda}(\beta)/a(\beta)$. If $m_{\mathrm{phys}}\in(0,\infty)$
and $a(\beta)\to0$, then $m_{\Lambda}(\beta)=m_{\mathrm{phys}}\,
a(\beta)(1+o(1))\to0$. A uniform lower bound
$m_{\Lambda}\ge c>0$ would give
$m_{\mathrm{phys}}=\lim m_{\Lambda}/a\ge c/\lim a=+\infty$,
contradicting $m_{\mathrm{phys}}<\infty$.
\end{proof}

\begin{theorem}[Attempted RG covariance and its failure]
\label{thm:rg_covariance_attempt}
Consider the following ``RG consistency'' argument for the
invariance of $h_{\mathrm{inv}}$:
\begin{quote}
\textit{Under RG blocking $R_{b}$: $m_{\Lambda}\to m_{\Lambda}/b$,
$D\to b(\beta'/\beta)D$. The Cheeger bound
$m_{\Lambda}\ge h_{\mathrm{RP}}^{2}/(2D)$ at blocked coupling
gives $h_{\mathrm{RP}}'{}^{2}\le h_{\mathrm{RP}}^{2}(\beta'/\beta)
(1+\delta_{\mathrm{gap}})$. By the ratio structure of $h_{\mathrm{RP}}=Q/\mu$,
the one-loop Jacobian cancels, giving $\delta_{\mathrm{gap}}=O(1/\beta^{2})$.
Hence $h_{\mathrm{inv}}'=h_{\mathrm{inv}}(1+O(1/\beta^{2}))$.}
\end{quote}
This argument is \textbf{invalid}: the Cheeger bound
$m_{\Lambda}\ge h_{\mathrm{RP}}^{2}/(2D)$ is an
\emph{inequality}, not an equality. The RG transformation of
$h_{\mathrm{RP}}$ is not constrained by the RG transformation of
$m_{\Lambda}$, because the ratio
$m_{\Lambda}/(h_{\mathrm{RP}}^{2}/(2D))\ge1$ can change under RG.
The ``consistency'' argument implicitly assumes the bound is
tight ($\delta_{\mathrm{gap}}=0$), which is false.
\end{theorem}

\begin{proof}
The error is in the step that derives
$h_{\mathrm{RP}}'{}^{2}\le h_{\mathrm{RP}}^{2}(\beta'/\beta)(1+\delta)$
from the gap transformation. The Cheeger bound at the blocked
coupling is $m_{\Lambda}'\ge h_{\mathrm{RP}}'{}^{2}/(2D')$, which
gives $h_{\mathrm{RP}}'{}^{2}\le 2D'm_{\Lambda}'$. Substituting
$m_{\Lambda}'=m_{\Lambda}/b$ and $D'=b(\beta'/\beta)D$:
\[
  h_{\mathrm{RP}}'{}^{2}\le 2b(\beta'/\beta)D\cdot m_{\Lambda}/b
  =2(\beta'/\beta)D\,m_{\Lambda}.
\]
But $m_{\Lambda}$ is the TRUE gap, not the Cheeger lower bound.
The Cheeger bound gives $m_{\Lambda}\ge h_{\mathrm{RP}}^{2}/(2D)$,
i.e.\ $2D\,m_{\Lambda}\ge h_{\mathrm{RP}}^{2}$, so
$2(\beta'/\beta)D\,m_{\Lambda}\ge(\beta'/\beta)h_{\mathrm{RP}}^{2}$.
This gives only
\[
  h_{\mathrm{RP}}'{}^{2}\le 2(\beta'/\beta)D\,m_{\Lambda},
\]
which is an \emph{upper} bound on $h_{\mathrm{RP}}'{}^{2}$ in
terms of the TRUE gap $m_{\Lambda}$, not in terms of
$h_{\mathrm{RP}}$. There is no constraint relating
$h_{\mathrm{RP}}'$ to $h_{\mathrm{RP}}$.

The Jacobian cancellation in the ratio $Q/\mu$ is a genuine
algebraic fact (see Lemma~\ref{lem:jacobian_cancel} below), but
it controls the transformation of the Cheeger \emph{constant}
$h_{\mathrm{RP}}=\inf_{A}Q(A)/\mu(A)$, not the gap bound
$m_{\Lambda}\ge h_{\mathrm{RP}}^{2}/(2D)$. The two are related
by an inequality that is not tight.
\end{proof}

\begin{lemma}[Jacobian cancellation in the Cheeger ratio]
\label{lem:jacobian_cancel}
The Cheeger constant $h_{\mathrm{RP}}=\inf_{A}Q(A)/\mu(A)$ is a
ratio of the boundary measure $Q$ to the volume measure $\mu$.
Under the block-spin RG $R_{b}$, both $Q$ and $\mu$ transform by
the same RG Jacobian $J=1+c_{1}g^{2}+O(g^{4})$. In the ratio
$Q/\mu$, the leading $O(g^{2})=O(1/\beta)$ Jacobian factors
cancel:
\begin{equation}
\label{eq:jacobian_cancel}
  \frac{Q'}{\mu'}=\frac{Q\cdot J}{\mu\cdot J}\cdot(1+O(g^{4}))
  =\frac{Q}{\mu}\cdot(1+O(1/\beta^{2})).
\end{equation}
This cancellation is verified explicitly for the single-plaquette
measure by the Peter--Weyl decomposition: the character expansion
$e^{\beta\chi(U)}=\sum_{\pi}d_{\pi}\frac{I_{\pi}(\beta)}{I_{0}(\beta)}
\chi_{\pi}(U)$ has the $\beta$-dependent normalization
$I_{0}(\beta)$ appearing in \emph{both} the boundary and volume
contributions to $h_{\mathrm{plaq}}$, hence cancelling in the
ratio.
\end{lemma}

\begin{remark}[What the Jacobian cancellation does and does not imply]
\label{rem:jacobian_scope}
The Jacobian cancellation (Lemma~\ref{lem:jacobian_cancel}) is a
genuine algebraic fact about the Cheeger \emph{ratio}
$h_{\mathrm{RP}}=Q/\mu$. It implies that $h_{\mathrm{RP}}$
transforms under RG as
\begin{equation}
\label{eq:h_rp_transform}
  h_{\mathrm{RP}}(R_{b}(\beta))
  =h_{\mathrm{RP}}(\beta)\cdot(\text{geometric factor})\cdot
  (1+O(1/\beta^{2})),
\end{equation}
where the ``geometric factor'' accounts for the change in the
isoperimetric profile under blocking (and is $O(1)$ for the OS
temporal Cheeger constant). However, this does \emph{not} imply
$h_{\mathrm{inv}}=h_{\mathrm{RP}}/\sqrt{\beta}$ is invariant,
because the geometric factor may itself depend on $\beta$ in a
way that cancels the $\sqrt{\beta}$ normalization. In fact, by
Lemma~\ref{lem:uniform_gap_contradicts_AF}, $h_{\mathrm{inv}}$
\emph{cannot} be bounded below (else the uniform gap would
contradict AF), so the geometric factor must drive
$h_{\mathrm{inv}}\to0$ at the AF rate.
\end{remark}

\begin{theorem}[Correct AF scaling of $h_{\mathrm{inv}}$]
\label{thm:h_inv_af_scaling}
If the Clay conjecture holds (finite $m_{\mathrm{phys}}>0$ with
AF scaling), then the RG-covariant Cheeger constant satisfies
\begin{equation}
\label{eq:h_inv_af_scaling}
  h_{\mathrm{inv}}(\beta)\;\sim\;\sqrt{\frac{a(\beta)}{\beta}}
  \;\sim\;\frac{e^{-b_{0}\beta/2}}{\sqrt{\beta}}\,(b_{0}\beta)^{b_{1}/(4b_{0})}
  \;\xrightarrow{\beta\to\infty}\;0.
\end{equation}
\end{theorem}

\begin{proof}
By the Cheeger bound $m_{\Lambda}\ge h_{\mathrm{RP}}^{2}/(2D)$
with $D=\beta N|P_{t}|$, and $m_{\Lambda}\sim m_{\mathrm{phys}}\,
a(\beta)$ (Clay + AF):
\[
  m_{\mathrm{phys}}\,a(\beta)\;\gtrsim\;\frac{h_{\mathrm{RP}}^{2}}{2\beta N|P_{t}|},
\]
so $h_{\mathrm{RP}}\lesssim\sqrt{2m_{\mathrm{phys}}\,N|P_{t}|\,\beta\,a(\beta)}$,
and $h_{\mathrm{inv}}=h_{\mathrm{RP}}/\sqrt{\beta}\lesssim
\sqrt{2m_{\mathrm{phys}}\,N|P_{t}|\,a(\beta)}\sim\sqrt{a(\beta)}$.
For the lower bound, by optimality of Lawler--Sokal,
$h_{\mathrm{RP}}\gtrsim\sqrt{m_{\Lambda}\cdot D}\sim\sqrt{a(\beta)\cdot\beta}$,
so $h_{\mathrm{inv}}\gtrsim\sqrt{a(\beta)}$. Combining:
$h_{\mathrm{inv}}\sim\sqrt{a(\beta)}\sim e^{-b_{0}\beta/2}\cdots$.
\end{proof}

\begin{remark}[Honest status: the obstruction is the geometric factor]
\label{rem:rg_honest_revised}
The analysis of this subsection reveals the \emph{true}
obstruction to proving Clay via the Cheeger approach:
\begin{enumerate}
\item The Jacobian cancellation (Lemma~\ref{lem:jacobian_cancel})
is valid and gives $h_{\mathrm{RP}}$ a controlled RG
transformation \eqref{eq:h_rp_transform} at two-loop order.
\item However, the ``geometric factor'' in
\eqref{eq:h_rp_transform} is not controlled by the Jacobian
cancellation. It encodes the change in the isoperimetric profile
of the OS measure under blocking.
\item By Theorem~\ref{thm:h_inv_af_scaling}, if Clay holds then
$h_{\mathrm{inv}}\to0$ at the AF rate $\sqrt{a(\beta)}$. Proving
this AF-rate decay \emph{from below} (i.e., showing
$h_{\mathrm{inv}}\ge c\sqrt{a(\beta)}$ for some $c>0$) is
equivalent to Clay.
\item The ``RG consistency'' argument
(Theorem~\ref{thm:rg_covariance_attempt}) is invalid because it
treats the Cheeger inequality as an equality.
\end{enumerate}
The residual obstruction is therefore: \emph{control the
geometric factor in \eqref{eq:h_rp_transform} to show it
contributes exactly the AF scaling $\sqrt{a(\beta)/\beta}$ to
$h_{\mathrm{inv}}$.} This is a concrete, well-defined analytic
problem on the OS isoperimetric profile, equivalent to Clay.
It is \emph{not} solved in this appendix.
\end{remark}

\begin{remark}[Comparison of the three Cheeger approaches]
\label{rem:three_approaches}
The three renormalized Cheeger constants introduced across
app217--app219 capture different facets of the Clay obstruction:
\begin{center}
\begin{tabular}{lll}
\hline
Approach & Constant & Target \\
\hline
CSL (app217) & $\|\mathcal{J}_{t_{*}}\|_{\mathrm{op}}$ & bounded after RG \\
RP-Cheeger (app218) & $h_{\mathrm{RP}}$ & $\inf_{\beta}h_{\mathrm{RP}}^{2}/(2D)>0$ \\
AF-renormalized (app219) & $h_{\mathrm{ren}}=h_{\mathrm{RP}}/\sqrt{a}$ & $\inf h_{\mathrm{ren}}>0$ \\
RG-covariant (app219) & $h_{\mathrm{inv}}=h_{\mathrm{RP}}/\sqrt{\beta}$ & $\ge c\sqrt{a(\beta)}$ \\
\hline
\end{tabular}
\end{center}
The app218 and app219 AF-renormalized approaches target
\emph{uniform} lower bounds (which are either false under AF or
require the circular Assumption~\ref{ass:fv_cluster}). The
app219 RG-covariant approach correctly identifies the target as
$h_{\mathrm{inv}}\ge c\sqrt{a(\beta)}$ (AF-rate decay), but
proving this requires control of the geometric factor in the RG
transformation of $h_{\mathrm{RP}}$, which is the residual
obstruction.
\end{remark}

\endinput
